%% ch11 --- Testy: pytest dla tych, którzy znają xUnit, i szkielet trace-assert
%%
%% Napisane we wrześniu 2026 pod pytest \pytestver, pytest-asyncio
%% \pytestasyncver, hypothesis \hypothesisver i coverage \coveragever.
%% Każdy listing to plik w code/ch11/, który uruchamia własne testy, gdy
%% odpalisz go jako skrypt, i CI uruchamia je wszystkie. Trzy transkrypty
%% generuje code/measure/transcripts.py, dwa z nich z celowo NIEUDANYCH
%% uruchomień plików trap_*.py, których zestaw testów nigdy nie zbiera.

\chapter{Testy: pytest dla tych, którzy znają xUnit, i szkielet trace-assert}
\label{ch:testing}

\noindent
pytest to xUnit z usuniętą klasą i z jawnie wypisanym grafem fixture'ów. To
całe mapowanie i prawie wszystko inne z niego wynika: nie ma typu bazowego
do odziedziczenia, atrybutu do przypięcia ani konstruktora do przeciążenia,
więc to, co one niosły \dash{} setup, teardown, współdzielony stan,
przypadki sterowane danymi \dash{} musiało trafić gdzieś indziej, a tym
gdzieś jest graf nazwanych funkcji, o które test prosi, wypisując ich nazwy.

Części, które się przenoszą, to te, których się spodziewasz. Ta, która się
nie przenosi, jest źródłem incydentów i nie jest nią graf fixture'ów: jest
nią podmiana zależności, przy której test w Pythonie potrafi przejść z
powodu niemającego nic wspólnego z testowanym kodem. To
\S\ref{sec:testing-patching} i to jest powód, dla którego ten rozdział
istnieje w książce dla ludzi, którzy już umieją testować.

\section{Co pytest zabrał}
\label{sec:testing-discovery}
\index{pytest!wykrywanie testów}

Moduł, którego nazwa zaczyna się od \code{test\_}, zawierający funkcje o
nazwach zaczynających się od \code{test\_}, jest zestawem testów. Nie ma
\code{[Fact]}, nie ma \code{TestClass} i nie ma czego rejestrować.

\pyregion{code/ch11/test_pricing.py}{plain}{Test to funkcja w pliku i tyle}{lst:ch11-plain}

\noindent
Każdy listing w tym rozdziale to plik, który uruchomisz na dwa sposoby
\dash{} pod pytest albo jako skrypt, dzięki trzem linijkom na jego końcu:

\pyregion{code/ch11/test_pricing.py}{runner}{Trzy linijki, dzięki którym plik z testami uruchamia się sam}{lst:ch11-runner}

\begin{csbox}
\code{[Fact]} staje się konwencją nazewniczą; \code{[Theory]} z
\code{[InlineData]} staje się \api{parametrize}
(\S\ref{sec:testing-parametrize}); konstruktor i \code{IDisposable} stają
się fixture z \code{yield} w środku; \code{IClassFixture<T>} staje się
zakresem. Wzorzec wykrywania rządzi \emph{szukaniem}: podaj plik w linii
poleceń, a pytest go zbierze niezależnie od nazwy \dash{} tak uruchamiane są
dwa celowo nieudane pliki z tego rozdziału.
\end{csbox}

\noindent
To, co zastępuje \code{Assert.Equal}, jest dziwniejsze i warto się przy tym
zatrzymać w pierwszej kolejności. Piszesz goły \code{assert} \dash{} słowo
kluczowe języka, to samo, które znika pod \code{-O} \dash{} a pytest
przepisuje je przy imporcie tak, żeby porażka sama się tłumaczyła. Oto co
wypisują dwie porażki z \code{code/ch11/trap\_assert.py}, wygenerowane przez
uruchomienie tego pliku, a nie zapamiętane:

\transcript{ch11-assert-rewrite}

\noindent
Nikt nie napisał komunikatu. pytest odtworzył podwyrażenie, które dało tę
wartość, a dla kolekcji znalazł różniący się indeks. Biblioteka asercji
istnieje po to, żeby ci to kupić, a tutaj słowo kluczowe języka i biblioteka
są tym samym.

\begin{warning}
Uruchom ten sam plik na CI, a ostatnia linia będzie inna:
zamiast \code{Use -v to get more diff} dostaniesz cały diff z
\api{difflib}, a długie wyjście przestanie być przycinane.
\api{running\_on\_ci} w \code{\_pytest.compat} jest prawdziwe, gdy
\code{CI} albo \code{BUILD\_NUMBER} jest ustawione i niepuste, a czytają
je dwa miejsca \dash{} porównanie sekwencji i przycinanie. To świadoma
uprzejmość, bo na runnerze nie powtórzysz uruchomienia z \code{-v}.
Znaczy też, że raport nie jest odtwarzalny między twoją maszyną a
serwerem, dopóki nie powiesz, który chcesz, więc skrypt generujący każdy
transkrypt w tej książce czyści obie zmienne.
\end{warning}

\begin{trapbox}
\code{assert} to nie jest ta instrukcja z buildów debugowych, o której
myślisz. To jest biblioteka asercji, bo pytest przepisuje bajtkod każdego
modułu, który zbierze. Konsekwencja jest jednak realna: uruchom interpreter
z \code{-O}, a każdy \code{assert} w twojej \emph{aplikacji} zniknie \dash{}
więc \code{assert} jest do testów i do niezmienników, które możesz stracić,
nigdy do walidacji danych wejściowych. Dodatek~\ref{app:traps} ma ten wpis.
\end{trapbox}

\noindent
Oczekiwane wyjątki to menedżer kontekstu zamiast atrybutu, a argument
\code{match} to wyrażenie regularne dopasowywane do komunikatu:

\pyregion{code/ch11/test_pricing.py}{raises}{Assert.Throws jako menedżer kontekstu}{lst:ch11-raises}

\section{Fixture: konstruktor zamieniony w graf}
\label{sec:testing-fixtures}
\index{pytest!fixture}

Fixture to funkcja z dekoratorem, a test prosi o nią, wymieniając ją jako
parametr. Nic nie jest wstrzykiwane po typie i nie ma kontenera: nazwa
parametru \emph{jest} wyszukiwaniem.

\pyregion{code/ch11/test_pricing.py}{fixture}{O fixture prosi się po nazwie, nigdy po typie}{lst:ch11-fixture}

\noindent
Fixture sam może poprosić o fixture i tu słowo graf zaczyna zarabiać na
siebie \dash{} i tu różnica wobec konstruktora przestaje być kosmetyczna.
Konstruktor wykonuje się dla każdego testu w klasie, niezależnie od tego,
czy test tego chciał; fixture wykonuje się tylko dla testów, które go
nazwały, i tylko tak głęboko, jak go nazwały.

\mermaidfig{ch11-fixture-graph}{Test nazywa to, czego potrzebuje, pytest rozwiązuje każdą nazwę do fixture, a fixture może nazwać kolejny}{graf fixture'ów}

\pyregion{code/ch11/test_pricing.py}{graph}{Fixture proszący o fixture}{lst:ch11-graph}

\noindent
Teardown to \code{yield}, a nie druga metoda. Wszystko przed nim to setup,
wszystko po nim wykonuje się, gdy test się kończy \dash{} także wtedy, gdy
test nie przeszedł, co jest dokładnie tą gwarancją, którą daje \code{using},
i powodem, dla którego \code{IAsyncLifetime} ma dwie metody:

\pyregion{code/ch11/test_pricing.py}{teardown}{Setup, yield, teardown, w jednej funkcji}{lst:ch11-teardown}

\noindent
Jak długo fixture żyje, decyduje jeden argument, a cztery wartości mapują
się na rzeczy, które już uruchamiasz:

\begin{center}
\small
\begin{tabularx}{\linewidth}{@{}lX@{}}
\toprule
\textbf{\code{scope=}} & \textbf{Czemu odpowiada} \\
\midrule
\code{"function"} & Konstruktor: nowy na każdy test. Domyślny i ten, przy
  którym warto zostać, dopóki coś nie zaboli. \\
\code{"class"} & \code{IClassFixture<T>}: współdzielony przez testy w jednej
  klasie, w rzadkim pliku pytest, który w ogóle ma klasę. \\
\code{"module"} & Współdzielony przez jeden plik. \\
\code{"session"} & \code{ICollectionFixture<T>} nad całością: budowany raz
  na cały przebieg. Kontener z bazą mieszka tutaj; mutowalny stan nie
  może. \\
\bottomrule
\end{tabularx}
\end{center}

\noindent
Fixture, którego potrzebuje kilka plików, trafia do \code{conftest.py} obok
nich. Ten plik nie jest przez nic importowany i nigdzie nie jest wymieniany:
pytest znajduje go, idąc w górę od każdego pliku z testami, a każdy fixture
w nim jest dostępny dla wszystkiego poniżej. To najbliższa rzecz, jaką
pytest ma do korzenia kompozycji, a Rozdział~\ref{ch:imports} dowodził, że
korzeń kompozycji to funkcja, a nie kontener.

\begin{trapbox}
Nie ma klasy testowej, na której można by powiesić fixture, a sięgnięcie po
nią jest pierwszym odruchem inżyniera .NET. Plik pytest z klasą to zwykle
plik, którego autor wciąż pisał xUnit; klasa kupuje grupowanie i
\code{scope="class"}, a kosztuje poziom wcięcia i \code{self}, którego nikt
nie czyta. Zacznij bez niej.
\end{trapbox}

\begin{exercise}{e11_01_fixture}{Usuń duplikację fixture'em}
Dwa testy w pliku startowym budują ten sam koszyk, linijka po linijce.
Zastąp duplikację jednym fixture o nazwie \code{basket} i spraw, żeby oba
testy o niego prosiły. Test obok uruchamia pytest na twoim pliku, a potem
patrzy, jak go napisałeś: twoje testy muszą przechodzić, fixture musi
istnieć, a koszyk musi być budowany raz.
\end{exercise}

\section{Jeden test, cztery przypadki}
\label{sec:testing-parametrize}
\index{pytest!parametrize}

\api{parametrize} to \code{[Theory]} i \code{[InlineData]} w jednym
dekoratorze. Pierwszy argument nazywa parametry, drugi to lista przypadków,
a pytest generuje po jednym teście na przypadek, z identyfikatorem, po
którym można wybierać w linii poleceń.

\pyregion{code/ch11/test_pricing.py}{parametrize}{Cztery przypadki, jedna funkcja testowa}{lst:ch11-parametrize}

\begin{csbox}
Każdy przypadek jest w raporcie osobnym testem, dokładnie tak jak wiersz
\code{[InlineData]}: cztery przejścia zamiast jednego, a porażka nazywa
przypadek, który nie przeszedł. Tam gdzie \code{[MemberData]} wzięłoby
właściwość, \api{parametrize} bierze dowolny obiekt iterowalny, więc
składanie listy albo odczyt pliku przy imporcie działają bez ceremonii
\dash{} a nałożenie dwóch dekoratorów \api{parametrize} na jedną funkcję
daje iloczyn kartezjański.
\end{csbox}

\begin{exercise}{e11_02_parametrize}{Zwiń cztery testy w jeden}
Plik startowy ma cztery testy, które różnią się tylko liczbami. Zrób z nich
jeden test sparametryzowany, zachowując wszystkie cztery przypadki.
Sprawdzenie wymaga dokładnie jednej funkcji testowej w pliku, dekoratora
parametrize i czterech testów, które nadal przechodzą.
\end{exercise}

\section{Podmiana zależności i pułapka}
\label{sec:testing-patching}
\index{pytest!monkeypatch}

Oto jedna rzecz w tym rozdziale, która będzie cię kosztować popołudnie.

Python ma dwa narzędzia do podmiany. \api{monkeypatch} to wbudowany fixture,
który ustawia atrybut i cofa to przy teardownie; \pkg{unittest.mock} to Moq
z biblioteki standardowej, z nagrywaniem i weryfikacją. Oba biorą rzecz do
podmiany jako \emph{napis}, a wybranie tego napisu źle daje ci test, który
przechodzi z niewłaściwego powodu albo nie przechodzi bez widocznej
przyczyny.

Testowany moduł robi najzwyklejszą rzecz w Pythonie:

\begin{python}
# code/ch11/invoice.py
from rates import vat_rate_percent

def gross_bound(net_pence: int) -> int:
    return net_pence + net_pence * vat_rate_percent() // 100
\end{python}

\noindent
Podmieniasz więc stawkę. Jest zdefiniowana w \code{rates}, więc to nazywasz.
Tym testem jest \code{trap\_patch.py}, trzymany w repozytorium po to, żeby
jego porażka była prawdziwa:

\transcript{ch11-patch-trap}

\noindent
Podmiana się wykonała. Stawka się nie zmieniła. Nic nie ostrzegło.

Rozdział~\ref{ch:imports} ma powód i mieści się on w jednym zdaniu: moduł to
obiekt, który wykonuje się raz. \code{from rates import vat\_rate\_percent}
wykonało się przy imporcie i związało \emph{obiekt} funkcji z przestrzenią
nazw modułu \code{invoice}. Są teraz dwie nazwy dla jednej funkcji, a
przepięcie jednej z nich zostawia drugą tam, gdzie wskazywała zawsze.

\mermaidfig{ch11-patch-target}{Jeden import, dwie nazwy dla jednego obiektu funkcji i tylko jedna z nich jest tą, którą czyta wywołanie}{gdzie musi trafić podmiana}

\begin{trapbox}
Podmieniaj tam, gdzie nazwa jest \emph{wyszukiwana}, a nie tam, gdzie jest
zdefiniowana. Która to nazwa, zdecydował import: \code{from x import y}
wiąże w module importującym, więc podmieniaj \code{importer.y};
\code{import x} i potem \code{x.y()} wyszukuje atrybut w chwili wywołania,
więc podmieniaj \code{x.y}. Reguła nie brzmi \enquote{zawsze podmieniaj
importera} \dash{} brzmi \enquote{podmieniaj nazwę, którą rozwiązuje
wołający kod}. Dodatek~\ref{app:traps} ma ten wpis i jest to wpis, który
najbardziej warto znać na pamięć.
\end{trapbox}

\noindent
Obie połowy są asercjami w \code{code/ch11/test\_patching.py}, więc build
tej książki nie przejdzie w dniu, w którym którakolwiek przestanie być
prawdziwa. Właściwy cel:

\pyregion{code/ch11/test_patching.py}{right}{Podmiana nazwy, którą czyta wywołanie}{lst:ch11-right}

\noindent
A ten sam zły cel jest właściwym celem dla drugiego zapisu, i to właśnie
czyni z tego regułę o importach, a nie regułę o podmienianiu:

\pyregion{code/ch11/test_patching.py}{qualified}{Wywołanie kwalifikowane rozwiązuje nazwę w chwili wywołania, więc celem jest definicja}{lst:ch11-qualified}

\noindent
\pkg{unittest.mock} to druga połowa. \code{patch} jako menedżer kontekstu
jest zakresem, a rodzina \code{assert\_called} to \code{Verify}:

\pyregion{code/ch11/test_patching.py}{mock}{unittest.mock, czyli Moq z mniejszą liczbą lambd}{lst:ch11-mock}

\begin{warning}
Goły \api{Mock} odpowiada na każdy atrybut i przyjmuje każde wywołanie, więc
\code{rate.assert\_called\_onse\_with()} \dash{} zwróć uwagę na literówkę
\dash{} sam jest atrybutem mocka, zwraca mocka i przechodzi. \api{Mock} ma
argument \code{spec}, a \code{patch} ma \code{autospec=True} właśnie na to;
użyj jednego z nich wszędzie tam, gdzie dubler zastępuje coś o sygnaturze
wartej uszanowania.
\end{warning}

\begin{exercise}{e11_03_patch_target}{Nazwij cel, który czyta wywołanie}
Plik startowy trzyma jeden napis: cel podmiany, który prawie każdy pisze
najpierw, a który nie robi zupełnie nic. Zmień go na nazwę, którą
\code{gross\_bound} rozwiązuje w chwili wywołania. Test podmienia to, co
nazywa twój napis, a potem sprawdza, czy podmiana dotarła do wywołania.
\end{exercise}

\section{Fixture'y, których nie napisałeś}
\label{sec:testing-builtin}
\index{pytest!wbudowane fixture'y}

pytest dostarcza fixture'y do trzech rzeczy, które test najczęściej musi
przygotować ręcznie. \api{tmp\_path} to \api{Path} do pustego katalogu,
unikalnego dla testu i usuwanego po nim:

\pyregion{code/ch11/test_builtins.py}{tmp_path}{Katalog roboczy, którego nikt nie musi sprzątać}{lst:ch11-tmppath}

\noindent
\api{caplog} przechwytuje \emph{rekordy} logu, a nie sformatowane linie, i o
to właśnie chodzi: asercja na sformatowanym napisie wiąże test z
formatterem, a formatter jest tą częścią, która najpewniej się zmieni i
najmniej warta jest testowania.

\pyregion{code/ch11/test_builtins.py}{caplog}{Asercja na polach rekordu, a nie na wyrenderowanej linii}{lst:ch11-caplog}

\noindent
\api{capsys} odczytuje to, co zostało wypisane, a \code{readouterr()} zwraca
oba strumienie i je czyści.

\pyregion{code/ch11/test_builtins.py}{capsys}{Odczytanie stdout}{lst:ch11-capsys}

\noindent
Markery to \code{[Trait]}, a wybieranie to \code{-m} dla markerów i
\code{-k} dla fragmentu nazwy testu. Marker trzeba zarejestrować w
\code{pyproject.toml}, inaczej pytest ostrzega \dash{} niezarejestrowany
marker to prawie zawsze literówka, a marker z literówką to test, który po
cichu nigdy nie uruchamia się pod filtrem, który miał go wybrać.

\pyregion{code/ch11/test_builtins.py}{markers}{Marker i powód, dla którego trzeba go zadeklarować}{lst:ch11-markers}

\begin{exercise}{e11_04_tmp_path}{Niech katalog należy do testu}
Zaimplementuj \code{write\_report}, które bierze katalog i zapisuje w nim
jeden plik. Nie może tworzyć katalogu, sprzątać go ani wybierać jego nazwy:
test podaje \api{tmp\_path}, a funkcji, która sama decyduje, dokąd trafia
jej wynik, nie da się przetestować bez dotykania maszyny, na której działa.
\end{exercise}

\section{Testy asynchroniczne, testy własności i co mierzy coverage}
\label{sec:testing-async}
\index{pytest!testy asynchroniczne}

Test asynchroniczny to test z \code{async} z przodu. Nie ma opakowania, nie
ma \code{.Result} i nie ma \code{GetAwaiter}, bo \pkg{pytest-asyncio}
uruchamia korutynę na pętli, którą sam zarządza.

\pyregion{code/ch11/test_async_property.py}{async}{Test asynchroniczny, z pętlą obsłużoną za ciebie}{lst:ch11-async}

\begin{warning}
Działa to, bo \code{code/pyproject.toml} ustawia
\code{asyncio\_mode = "auto"}. W domyślnym trybie strict nieoznaczony test
asynchroniczny nie zostaje uznany za nieudany \dash{} zostaje
\emph{pominięty}, z ostrzeżeniem. Zestaw testów, który raportuje zieleń, nie
uruchomiwszy żadnego ze swoich testów asynchronicznych, to najgorszy możliwy
tu wynik, więc ustaw tryb raz, w pliku projektu, i nigdy per test.
Rozdział~\ref{ch:asyncio} to reszta tej historii.
\end{warning}

\noindent
hypothesis to FsCheck: podaj własność, a biblioteka znajdzie przypadek,
który ją łamie. Dodatkowo zwęża porażkę do najmniejszej postaci i zapamiętuje
ją, więc przypadek znaleziony raz jest odtwarzany przy każdym późniejszym
uruchomieniu, zamiast czekać, aż to samo losowanie wypadnie ponownie.

\pyregion{code/ch11/test_async_property.py}{hypothesis}{Własność, a wejścia wymyśla biblioteka}{lst:ch11-hypothesis}

\noindent
Coverage na koniec i krótko, bo ta liczba jest warta mniej niż jej reputacja.
\pkg{coverage} to coverlet: zapisuje, które linie się wykonały. Oto co mówi o
\code{code/ch11/pricing.py} pod testem, który woła każdą funkcję w nim i
sprawdza w gruncie rzeczy nic:

\transcript{ch11-coverage}

\noindent
Moduł jest pokryty w całości, w jedynym sensie, o jaki narzędziu chodzi.
Każda linia się wykonała; żadna z asercji w
\code{code/ch11/trap\_coverage.py} nie mówi, co którakolwiek z nich miała
zwrócić, więc moduł mógłby zwracać dowolną liczbę całkowitą, a raport
byłby ten sam. Coverage znajduje kod, do którego nie dociera żaden test, i
jest to rzecz realna i warta wiedzy. Nigdy nie potrafił powiedzieć, czy kod,
do którego dotarł, został sprawdzony, a cel wyrażony w procentach kupuje
testy, które wykonują, a nie testy, które asertują.

\section{trace-assert, etap 01}
\label{sec:testing-traceassert}
\index{trace-assert!etap 01}

\begin{projectbox}
Etap 01: model śladu, fixture \code{trace} i pierwsze dwie asercje, w
\code{code/src/trace\_assert/}. Etap 02 zapisuje do niego prawdziwe
wywołanie modelu w Rozdziale~\ref{ch:aitoolkit};
Rozdział~\ref{ch:traceassert} portuje resztę zestawu asercji, pakuje go i
porównuje trzy porty.
\end{projectbox}

\noindent
Projekt przewodni to port pierwszej warstwy agent-eval-bench:
deterministyczne asercje o tym, co agent \emph{zrobił}, wyrażone jako
funkcje pytest. Projekt tego rozwiązania nie należy do tej książki, a port
idzie za oryginałem, zamiast wymyślać go od nowa \dash{} trzy porty jednego
projektu to jedyne porównanie warte zrobienia w
Rozdziale~\ref{ch:traceassert}.

Dwa obiekty, a podział między nimi to całość etapu 01. \api{Recorder} jest
mutowalny i należy do przebiegu; \api{Trace} jest niezmienny i należy do
asercji, bo asercja, która może edytować dowód, nie jest asercją.

\mermaidfig{ch11-trace-assert}{Fixture daje testowi mutowalny Recorder; asercje czytają jego niezmienny zrzut}{kształt etapu 01 trace-assert}

\noindent
Fixture to cztery linijki w module wtyczki, a jedna linijka
\code{conftest.py} go rejestruje. Ta linijka staje się punktem wejścia w
metadanych samego pakietu w Rozdziale~\ref{ch:traceassert} i od tego
momentu wystarczy zainstalować trace-assert; etapowanie jest celowe, bo
punkt wejścia to pakowanie.

\pyfile{code/src/trace_assert/plugin.py}{Fixture trace jako moduł wtyczki}{lst:ch11-plugin}

\noindent
Potem dwie asercje. Schemat, z którego pochodzą, opisuje je jako parę: jedna
stwierdza, że narzędzie zostało wywołane, druga, że nie zostało, a ścieżka
odmowy potrzebuje obu \dash{} agent, który odmawia w tekście i mimo to woła
narzędzie, przechodzi tę pierwszą samodzielnie.

\pyfile{code/src/trace_assert/assertions.py}{Pierwsze dwie asercje i komunikaty, którymi się nie udają}{lst:ch11-assertions}

\noindent
Dwie rzeczy w środku należą do pytest, a nie do trace-assert. \code{\_\_%
tracebackhide\_\_} trzyma ten plik poza raportem porażki, więc czytelnik
widzi linijkę w \emph{swoim} teście i komunikat, a nie linijkę wewnątrz
biblioteki, która zgłosiła wyjątek. A podanie obu ograniczeń podnosi
\api{ValueError}, zamiast po cichu rozstrzygać: schemat oryginału też
odrzuca tę kombinację i podaje powód \dash{} jedno z nich zostałoby
zignorowane, a ty dalej wierzyłbyś w ograniczenie, którego nic nigdy nie
sprawdziło.

\begin{exercise}{e11_05_trace}{Zapewnij odmowę i nieobecność}
Agent od zwrotów może wypłacać do limitu, a powyżej niego musi eskalować.
Zaimplementuj \code{handle\_refund} tak, żeby zapisywał, co zrobił, a potem
przeczytaj test: stwierdza on wywołanie, które miało nastąpić, \emph{oraz}
to, które nastąpić nie może. Trafienie w tę drugą połowę jest tym
ćwiczeniem.
\end{exercise}

\noindent
To jest ten zestaw testów, który ta książka uruchamia od
Rozdziału~\ref{ch:runtime}: bez sieci, bez serwera bazy danych i bez modelu,
i to właśnie pozwala mu chodzić przy każdym pushu. Ocena tego, czy
\emph{odpowiedź} agenta była dobra, to inny problem i należy do tomu o
LangChain. Tutaj należy warstwa pod nią \dash{} ta deterministyczna, a więc
ta, która potrafi zablokować merge.
